\documentclass[10pt,a4paper]{amsart}
\usepackage[margin=19mm]{geometry}
\usepackage{amsmath,amssymb,mathtools}
\usepackage[scr=rsfs,cal=euler]{mathalfa}
\usepackage{xfrac}
\usepackage[unicode,hidelinks,bookmarks=false]{hyperref}
\newcommand{\ie}{i.e.\ }
\newcommand{\cf}{cf.\ }
\newcommand{\resp}{respectively\ }
\newcommand{\Subsection}{\subsection}
\providecommand{\nobreakdash}{\nobreak\hbox{-}}
\setlength{\emergencystretch}{2em}
\multlinegap0pt
\usepackage{bm}
\usepackage{bbm}
\usepackage{dsfont}
\usepackage{xspace}
\usepackage{cancel}
\usepackage{mathtools}
\usepackage{amsthm}
\makeatletter
\@ifundefined{theorem}{\newtheorem{theorem}{Theorem}}{}
\@ifundefined{corollary}{\newtheorem{corollary}[theorem]{Corollary}}{}
\@ifundefined{proposition}{\newtheorem{proposition}[theorem]{Proposition}}{}
\makeatother
\def\Z{\mathbb Z}
\def\aut#1{\mathrm{Aut}(#1)}
\title[Latin quandles of size $16p$]{Latin quandles of size $16p$}
\author{Marco Bonatto, Filippo Spaggiari}
\date{Source: arXiv:2509.22145v1 (CC BY 4.0)}
\begin{document}
\maketitle

\noindent\emph{Source and licence.} Latin quandles of size $16p$. Version arXiv:2509.22145v1 (CC BY 4.0), distributed under the Creative Commons Attribution 4.0 International licence, as recorded in the arXiv licence metadata for this exact version and shown on the abstract page https://arxiv.org/abs/2509.22145v1. Full rights notice: licenses/arxiv-2509.22145-CC-BY-4.0.txt.\par\smallskip

\noindent\textit{Editorial scope.} Counted unit: the Proposition whose source label is \texttt{chain 4 then p=3} in the article (source0.tex lines 3424--3426, source label \texttt{chain 4 then p=3}), delivered together with the article's own complete original proof of it (source0.tex lines 3428--3451, one \texttt{proof} environment of 24 lines). No mathematics is written, repaired or re-derived: statement and proof are the article's own text line for line. Required context retained uncounted: aut of G\_{M,N (lines 3339--3346); def A,B (lines 3383--3394). Every definition, display and label of those lines is retained unchanged. Omitted from this package and not redistributed: 1--3338, 3347--3382, 3395--3423, 3427--3427, 3452--3967. Nothing inside a retained excerpt is omitted or altered: every retained body is delivered exactly as the article's lines are written, so no omission receipt of any kind accompanies this package. Citations: the retained excerpts carry no citation command at all, so no bibliography entry of the article is transcribed and no reference is redistributed. Reference adaptation: none was needed and none was made; every reference inside the delivered unit resolves inside it, so no unresolved reference or citation is produced. Printed slips kept literal: no printed word, symbol, hypothesis or number of the counted statement or of its proof was altered. Layout and class: the article's own class is \texttt{amsart} with packages including rotating, listings, amsthm, amssymb, url, hyperref, fullpage, MnSymbol, xy, tikz, subfig, caption, forest, color; the wrapper is the lane's pinned \texttt{amsart} editorial wrapper at 10pt on A4, which restates the theorem-like environments the retained text needs and adds the article's own macro definitions that the retained text needs (\texttt{\textbackslash Z}, \texttt{\textbackslash aut}), with extra wrapper packages (bm, bbm, dsfont, xspace, cancel, mathtools). The delivered numbering is the unit's own. Assets: no retained span contains a figure environment, a graphics-inclusion command, a PDF-page inclusion, a table or a TikZ picture; no image file is copied into this package, the package declares no asset dependency and no image file is redistributed with it. Licence: version arXiv:2509.22145v1 of this article is distributed under the Creative Commons Attribution 4.0 International licence, as recorded in the arXiv licence metadata for this exact version and shown on the abstract page https://arxiv.org/abs/2509.22145v1. A full rights notice is delivered at licenses/arxiv-2509.22145-CC-BY-4.0.txt. Characters: every retained excerpt is pure ASCII, so the delivered engine, XeLaTeX with its default Latin Modern text font, reports no missing character.
\begin{corollary}\label{aut of G_{M,N}}
    The automorphism group of $G_{M,N}(t)$ are the maps defined by setting:
\begin{align*}
        f(0,(1,0))=(w,(r,s)),\quad 
        f(0,(0,1))=(v,(h,k)),\quad f(u,(0,0))=(G(u),(0,0))
    \end{align*}
where $\{(r,s),(h,k)\}$ generates $\Z_4^2$, $G\in GL_t(p)$ such that $GMG^{-1}=A^rB^s$ and $GNG^{-1}=A^hB^k$.% In particular $|\aut{G_M}|=p^{2t}|N_{GL_t(\Z_p)}(\langle M,N\rangle)|$.
\end{corollary}

 \begin{align}\label{def A,B}
       A=\begin{bmatrix}
           a & 0 & 0\\
           0 & 1& 0\\
           0 & 0 & a^{-1}
       \end{bmatrix} , \quad 
       B=\begin{bmatrix}
           1 & 0 & 0\\
           0 & a& 0\\
           0 & 0 & a^{-1}
       \end{bmatrix} 
    \end{align}

\begin{proposition}\label{chain 4 then p=3}
    Let $G_{A,B}(3)$ where $A,B$ are defined as in \eqref{def A,B}. There is no $f\in \aut{G_{A,B}(3)}$ such that $|Fix(f)|=p^2$ and $Fix(f_{\Z_p^3})=1$. 
\end{proposition}

\begin{proof}
Let $G=G_{A,B}(t)$. According to Corollary \ref{aut of G_{M,N}} then $FAF^{-1}=A^r B^s$ and $FBF^{-1}=A^h B^k$ and the induced map $f_{\Z_p^3}$ defined as $(1,0)\mapsto (r,s)$, $(0,1)\mapsto (h,k)$ is an automorphism of $\Z_4^2$. In particular, $FAF^{-1},FBF^{-1}$ are diagonal matrices with eigenvalues $a,1,a^{-1}$. Accordingly $FAF^{-1},FBF^{-1}\in \{A^{\pm 1}, B^{\pm 1}, (AB^{-1})^{\pm 1}\}$. Moreover, $Fix(f_{\Z_p^3})=1$, so the only choices allowed are $FAF^{-1}=B^{-1}$, $FBF^{-1}=AB^{-1}$ or $FAF^{-1}=AB^{-1}, FBF^{-1}=A^{-1}$ and accordingly $F$ is 
\begin{align}\label{F}
    F=\begin{bmatrix}
        0 & x & 0\\
        y & 0 & 0\\
    0& 0 & z
    \end{bmatrix}
\end{align}
where $\det(F)=-xyz\neq 0$. On the other hand, the matrix $F$ has an eigenspace relative to $1$ of dimension $2$ and so $F$ is similar to a matrix of the form
\begin{align}\label{F'}
    F'=\begin{bmatrix}
        1 & 0 & t\\
        0 & 1 & u\\
    0& 0 & w
    \end{bmatrix}
\end{align}
for some $w\neq 0$.
The characteristic polynomials of $F$ and $F'$ coincide, so we have
\begin{align*}
  k^3- xyz=(k-1)^2(k-w)=k^3-(2+w)k^2+(2w+1)k-w.
\end{align*}
Accordingly $w=-2$ and so $Tr(F')=0=Tr(F)=z$. Hence $\det(F)=0$, contradiction.
\end{proof}

\end{document}
